\documentclass[11pt]{article}
\usepackage[margin=0.82in]{geometry}
\usepackage[T1]{fontenc}
\usepackage{lmodern,amsmath,amssymb,amsthm,microtype}
\usepackage[colorlinks=true,urlcolor=blue]{hyperref}
\newtheorem{theorem}{Theorem}
\newcommand{\Q}{\mathbb Q}
\newcommand{\Z}{\mathbb Z}
\title{Global nonvanishing after lower-period\\cancellation at weight five}
\author{}
\date{}
\begin{document}
\maketitle
\begin{abstract}
Two explicit third-order recurrence certificates establish the adjacent
coefficient signs needed to cancel $\zeta(3)$ at every even index.
A uniform recurrence cone proves positivity from three initial values;
the argument does not extrapolate finite sign checks. The certificates
are checked independently in $\Z[n,t]$. Lean verifies the cone induction,
its coefficient inequalities, and positivity of the specified recurrence
states. Their identification with the Laurent coefficients is a paper
proof here. The winning primitive arithmetic bound, and hence
irrationality of $\zeta(5)$, remain unproved.
\end{abstract}

\section{The coefficients and the claim}
For $a=0,1$ and $n\ge2$, put
\[
 R_n^{(a)}(t)=(n!)^3
 \frac{\prod_{v=1}^{n-a}(t-v)(t+n+v)}
      {\prod_{j=0}^n(t+j)^5}.
\]
Let $\phi_3$ sum all third-order partial-fraction coefficients in the integer
pole orbit, and set $v_n^{(a)}=\phi_3(R_n^{(a)})$. In the original normalization,
\[
 T_{n,m}(t)=\frac{(n!)^5}{(m!)^2}
 \frac{\prod_{v=1}^m(t-v)(t+n+v)}{\prod_{j=0}^n(t+j)^5},
 \qquad
 \alpha_{n,n}=v_n^{(0)},\quad \alpha_{n,n-1}=n^2v_n^{(1)}.
\]
For even $n$, $\alpha_{n,m}$ is the coefficient of $\zeta(3)$ in
$\mathcal A_{n,m}:=\sum_{k\ge1}T_{n,m}(k)$.
The sums converge absolutely by rational decay and are strictly positive:
their summands vanish for $k\le m$ and are positive for $k>m$.
Partial-fraction summation identifies $\alpha$ with the stated coefficient.
Reflection $T_{n,m}(-n-t)=-T_{n,m}(t)$ cancels the even zeta terms;
proper decay gives $\sum_jc_{j,1}=0$, leaving a rational constant and
only $\zeta(3),\zeta(5)$.

\begin{theorem}
For every even $n\ge2$,
\[
 (-1)^{n/2}\alpha_{n,n}>0,\qquad
 (-1)^{n/2+1}\alpha_{n,n-1}>0.
\]
In particular $\alpha_{n,n}\alpha_{n,n-1}<0$. Consequently
\[
 (-1)^{n/2+1}\bigl[
 \alpha_{n,n-1}\mathcal A_{n,n}
 -\alpha_{n,n}\mathcal A_{n,n-1}\bigr]>0.
\]
The last expression belongs to $\Q+\Q\zeta(5)$.
\end{theorem}

The proof is computer assisted through finite polynomial certificates.
The two accompanying JSON files specify all polynomials below by integer
coefficient lists in ascending powers; the standard-library verifier
checks the displayed identities by complete coefficient expansion.

\section{An exact index recurrence}
Fix $a$. The certificate gives $q_0,\ldots,q_3$ and
$X(n,t)=\sum_{k=0}^{30}x_k(n)t^k$. Define
\[
 U=t^6(t+2n-a+1),\qquad V=(t-n+a-1)(t+n+6)^6,
\]
\[
 H_\ell(n,t)=
 \prod_{i=1}^{2\ell}(t-n+a-i)
 \prod_{i=1}^{4\ell}(t+2n-a+i)
 \prod_{i=2\ell+1}^{6}(t+n+i)^6.
\]
The first checked identity, in $\Z[n,t]$, is
\begin{equation}
 U(n,t)X(n,t+1)-V(n,t)X(n,t)
           =\sum_{\ell=0}^3q_\ell(n)H_\ell(n,t).
 \label{eq:cert}
\end{equation}
The files also specify $P_0,\ldots,P_3$ and verify
\[
 q_\ell(n)\prod_{i=2\ell+1}^{6}(n+i)^3=g_a(n)P_\ell(n),\qquad
 g_0=\prod_{i=3}^6(n+i)^3,\quad g_1=\prod_{i=2}^6(n+i)^3.
\]
All $P_\ell$ have degree $25$ for $a=0$ and degree $36$ for $a=1$.

To interpret the certificate, put
\[
 D(n,t)=\prod_{i=1}^6(t+n+i)^6,\quad
 B(n,t)=\frac{R_n^{(a)}(t)}{D(n,t)},\quad G(n,t)=V(n,t)X(n,t)B(n,t).
\]
Direct cancellation gives the rational-function identities
\[
 \frac{B(n,t+1)}{B(n,t)}
 =\frac{t^6(t+2n-a+1)}{(t-n+a)(t+n+7)^6},
 \qquad
 B(n,t)H_\ell(n,t)=
 \frac{R_{n+2\ell}^{(a)}(t)}{\prod_{i=1}^{2\ell}(n+i)^3}.
\]
Thus \eqref{eq:cert}, multiplied by $B$, gives a rational telescoper
whose right side is a linear combination of the four shifted kernels.

The functional $\phi_3$ is translation invariant: replacing $G(t)$ by
$G(t+1)$ shifts each integer pole by one without changing its third-order
coefficient. All these rational functions have finitely many integer poles.
Applying $\phi_3$, clearing the factorial factors, and using $g_a(n)>0$
therefore proves
\begin{equation}
 \sum_{\ell=0}^3P_\ell(n)v_{n+2\ell}^{(a)}=0\qquad(n\ge2).
 \label{eq:rec}
\end{equation}
This does not assume numerical linear independence of zeta values:
the functional is applied before evaluating any period.

\section{A cone that preserves the sign}
The verifier checks that each coefficient of
\[
 -P_0(x+2),\quad P_1(x+2),\quad P_2(x+2),\quad P_3(x+2)
\]
is a nonnegative integer, with a positive constant term. The same is true of
\[
 \Omega(x)=10^6P_2(x+2)-10^9P_3(x+2)-1000P_1(x+2)+P_0(x+2).
\]
These are polynomial positivity certificates for all real $x\ge0$,
not sign evaluations at finitely many indices. For $n\ge2$, set
\[
 A=-P_0(n)>0,\quad B=P_1(n)>0,\quad C=P_2(n)>0,\quad D=P_3(n)>0.
\]
Then
\begin{equation}
 10^6 C>10^9D+1000B+A.
 \label{eq:cone}
\end{equation}

Let $u_k^{(a)}=(-1)^{k+1+a}v_{2k+2}^{(a)}$. Equation \eqref{eq:rec} becomes
\[
 D u_{k+3}=C u_{k+2}-B u_{k+1}-A u_k,
\]
where the coefficients are evaluated at $n=2k+2$. Suppose
$u_k>0$, $u_{k+1}\ge1000u_k$, and $u_{k+2}\ge1000u_{k+1}$.
Then \eqref{eq:cone} gives
\[
 D u_{k+3}\ge
 \left(C-\frac B{1000}-\frac A{10^6}\right)u_{k+2}
 \ge1000D u_{k+2}.
\]
The cone is therefore preserved.

Direct Laurent expansion at the first three indices gives
\[
 (u_0^{(0)},u_1^{(0)},u_2^{(0)})
 =\left(882,\frac{9149875}{2},\frac{225341478859}{5}\right),
\]
\[
 (u_0^{(1)},u_1^{(1)},u_2^{(1)})
 =\left(\frac{785}{8},\frac{16148335}{128},\frac{89419424171}{162}\right).
\]
Both triples satisfy the cone by exact rational comparisons.
Induction proves $u_k^{(0)},u_k^{(1)}>0$ for every $k\ge0$, establishing
the two asserted coefficient signs. The cancellation weights
$\alpha_{n,n-1}$ and $-\alpha_{n,n}$ have the common sign
$(-1)^{n/2+1}$; strict positivity of both $\mathcal A$ sums proves the
last inequality in the theorem. The $\zeta(3)$ coefficient cancels
algebraically. No separate all-index assertion about the nonzero
$\zeta(5)$ coefficient is needed or made.

\section{Verification scope and remaining arithmetic}
The files
\texttt{pade-sign-recurrence-0.json} and
\texttt{pade-sign-recurrence-1.json} contain the complete certificates.
Run
\begin{quote}
 \texttt{python scripts/verify\_sign\_recurrences.py}
\end{quote}
from the included Lean project. This verifier uses only integer and
rational arithmetic from Python's standard library. SymPy 1.14.0 is
needed only to rediscover the certificates, not to check them.

The Lean module \texttt{PadeSignRecurrence.lean} proves the cone step,
its induction, the coefficient inequalities using nonnegative Horner
coefficient lists, and positivity of the two explicitly defined recurrence
states. The step identifying those states with the actual Laurent-coefficient
sequences uses the paper argument from \eqref{eq:cert} to \eqref{eq:rec}.
That connection is not yet formalized in Lean; no complete Lean theorem
about the global Laurent-coefficient signs is claimed here.

Nonvanishing after lower-period cancellation is now established for this
family at every even index. The primitive arithmetic multiplier remains
\[
 K_n=\frac{d_n^5\sum_j|w_{n,j}|}{h_n},
 \qquad d_n=\operatorname{lcm}(1,\ldots,n),
\]
where the $w_{n,j}$ are integer cancellation weights obtained by clearing
the two rational weights, and $h_n$ is the full coefficient gcd.
The needed limiting comparison
between $K_n$ and the analytic decay is still unproved.
This note does not prove irrationality of $\zeta(5)$ or of all odd zeta values.
\end{document}
